% Curs 11 - Pagina 6 (13-XII-1994).
% Fragment pentru % Curs 11 - Pagina 6 (13-XII-1994).
% Fragment pentru % Curs 11 - Pagina 6 (13-XII-1994).
% Fragment pentru % Curs 11 - Pagina 6 (13-XII-1994).
% Fragment pentru \input{courses/pages/curs11-pagina6-ro}.
% Fara preambul si fara intrerupere fortata de pagina.
% Transcriere fidela; incertitudinile sunt mentionate in comentarii.

\subsection*{Curs 11 - Pagina 6}

% Continuarea demonstratiei teoremei de izomorfism, de la pagina 5.
\noindent $\underbrace{\bar a}_{aH}\in\ker f^*
\iff aH=aN\iff a\in H$.
% In primul rand al manuscrisului, aH=aN si a\in H sunt citiri
% probabile, dar compatibilitatea lor cu morfismul nu este clara.
\[
\ker f^*=\{x\mid (\ldots)\} = \frac{H}{N}.
\]
% Textul dintre bara verticala si acolada este partial barat/ilizibil.
% De asemenea, H/N este ceea ce pare scris; a nu se interpreta
% ca o corectie matematica a nucleului morfismului.

\bigskip
\noindent\underline{Aplica\c{t}ii:} \quad $(\mathbb Z,+,0)$.

\medskip
\noindent $H\leq\mathbb Z\ \Rightarrow\ (\exists!)\,u\in\mathbb N^*$
 astfel \^{i}nc\^{a}t $H=u\mathbb Z$.
% Manuscrisul pare sa foloseasca N^*; cazul H={0} cere u=0,
% deci enuntul, luat literal, nu acopera subgrupul trivial.

\bigskip
\noindent Fie $m,n\in\mathbb N^*$. Ar\u{a}t c\u{a}
\[
mn=[m,n]\cdot(m,n).
\]
\noindent % Notatii din manuscris: [m,n] = c.m.m.m.c.; (m,n) = c.m.m.d.c.
\[
m\mathbb Z+n\mathbb Z\leq\mathbb Z
\ \Longrightarrow\ (\exists!)\,d\in\mathbb N\ \text{a.\^{i}.}\
m\mathbb Z+n\mathbb Z=d\mathbb Z.
\]
\noindent Ar\u{a}t c\u{a} $d=(m,n)$.

\medskip
\noindent $(\exists)\,x,y\in\mathbb Z$ a.\^{i}.
\[
mx+ny=d
\quad\Longrightarrow\quad
\text{dac\u{a} } c\mid m\ \text{\c{s}i}\ c\mid n,
\ \text{atunci}\ c\mid d.
\]
\[
m\mathbb Z\subseteq d\mathbb Z,\qquad
m\in m\mathbb Z\ \Longrightarrow\ (\exists)\,q\in\mathbb Z
\ \text{a.\^{i}.}\ m=dq,\quad d\mid m.
\]
\noindent Analog $d\mid n$ $\Longrightarrow d=(m,n)$.

\bigskip
\noindent $m\mathbb Z\cap n\mathbb Z\leq\mathbb Z\ \Longrightarrow\
(\exists!)\,u\in\mathbb N$ a.\^{i}.
\[
m\mathbb Z\cap n\mathbb Z=u\mathbb Z.
\]
\noindent $u\in m\mathbb Z$ \c{s}i $u\in n\mathbb Z$
$\Longrightarrow m\mid u$ \c{s}i $n\mid u$.

\noindent Dac\u{a} $m\mid b$ \c{s}i $n\mid b$
$\stackrel{?}{\Longrightarrow} b\mid u$.
% In manuscris relatia finala pare b|u, dar randul urmator
% (b\in mZ si b\in nZ => b\in uZ) implica u|b. Verificare necesara.
\[
\underbrace{b\in m\mathbb Z}_{m\mid b},\quad
\underbrace{b\in n\mathbb Z}_{n\mid b}
\quad\Longrightarrow\quad b\in u\mathbb Z.
\]

\bigskip
\noindent T I de iz. $\Longrightarrow$
\[
\frac{m\mathbb Z+n\mathbb Z}{n\mathbb Z}
\cong \frac{m\mathbb Z}{m\mathbb Z\cap n\mathbb Z}.
\]
\[
\Longrightarrow\qquad
\underbrace{\frac{d\mathbb Z}{n\mathbb Z}}_{\cong\,\mathbb Z_{n_1}}
\cong
\underbrace{\frac{m\mathbb Z}{u\mathbb Z}}_{\cong\,\mathbb Z_{u_1}}.
\]
% Indicii n_1 si u_1 de la baza paginii sunt usor estompati.
\[
u=m u_1,\qquad n=d n_1.
\]
\[
u_1=n_1\quad\Longrightarrow\quad
\frac{u}{m}=\frac{n}{d}\quad\Longrightarrow\quad du=mn.
\]
.
% Fara preambul si fara intrerupere fortata de pagina.
% Transcriere fidela; incertitudinile sunt mentionate in comentarii.

\subsection*{Curs 11 - Pagina 6}

% Continuarea demonstratiei teoremei de izomorfism, de la pagina 5.
\noindent $\underbrace{\bar a}_{aH}\in\ker f^*
\iff aH=aN\iff a\in H$.
% In primul rand al manuscrisului, aH=aN si a\in H sunt citiri
% probabile, dar compatibilitatea lor cu morfismul nu este clara.
\[
\ker f^*=\{x\mid (\ldots)\} = \frac{H}{N}.
\]
% Textul dintre bara verticala si acolada este partial barat/ilizibil.
% De asemenea, H/N este ceea ce pare scris; a nu se interpreta
% ca o corectie matematica a nucleului morfismului.

\bigskip
\noindent\underline{Aplica\c{t}ii:} \quad $(\mathbb Z,+,0)$.

\medskip
\noindent $H\leq\mathbb Z\ \Rightarrow\ (\exists!)\,u\in\mathbb N^*$
 astfel \^{i}nc\^{a}t $H=u\mathbb Z$.
% Manuscrisul pare sa foloseasca N^*; cazul H={0} cere u=0,
% deci enuntul, luat literal, nu acopera subgrupul trivial.

\bigskip
\noindent Fie $m,n\in\mathbb N^*$. Ar\u{a}t c\u{a}
\[
mn=[m,n]\cdot(m,n).
\]
\noindent % Notatii din manuscris: [m,n] = c.m.m.m.c.; (m,n) = c.m.m.d.c.
\[
m\mathbb Z+n\mathbb Z\leq\mathbb Z
\ \Longrightarrow\ (\exists!)\,d\in\mathbb N\ \text{a.\^{i}.}\
m\mathbb Z+n\mathbb Z=d\mathbb Z.
\]
\noindent Ar\u{a}t c\u{a} $d=(m,n)$.

\medskip
\noindent $(\exists)\,x,y\in\mathbb Z$ a.\^{i}.
\[
mx+ny=d
\quad\Longrightarrow\quad
\text{dac\u{a} } c\mid m\ \text{\c{s}i}\ c\mid n,
\ \text{atunci}\ c\mid d.
\]
\[
m\mathbb Z\subseteq d\mathbb Z,\qquad
m\in m\mathbb Z\ \Longrightarrow\ (\exists)\,q\in\mathbb Z
\ \text{a.\^{i}.}\ m=dq,\quad d\mid m.
\]
\noindent Analog $d\mid n$ $\Longrightarrow d=(m,n)$.

\bigskip
\noindent $m\mathbb Z\cap n\mathbb Z\leq\mathbb Z\ \Longrightarrow\
(\exists!)\,u\in\mathbb N$ a.\^{i}.
\[
m\mathbb Z\cap n\mathbb Z=u\mathbb Z.
\]
\noindent $u\in m\mathbb Z$ \c{s}i $u\in n\mathbb Z$
$\Longrightarrow m\mid u$ \c{s}i $n\mid u$.

\noindent Dac\u{a} $m\mid b$ \c{s}i $n\mid b$
$\stackrel{?}{\Longrightarrow} b\mid u$.
% In manuscris relatia finala pare b|u, dar randul urmator
% (b\in mZ si b\in nZ => b\in uZ) implica u|b. Verificare necesara.
\[
\underbrace{b\in m\mathbb Z}_{m\mid b},\quad
\underbrace{b\in n\mathbb Z}_{n\mid b}
\quad\Longrightarrow\quad b\in u\mathbb Z.
\]

\bigskip
\noindent T I de iz. $\Longrightarrow$
\[
\frac{m\mathbb Z+n\mathbb Z}{n\mathbb Z}
\cong \frac{m\mathbb Z}{m\mathbb Z\cap n\mathbb Z}.
\]
\[
\Longrightarrow\qquad
\underbrace{\frac{d\mathbb Z}{n\mathbb Z}}_{\cong\,\mathbb Z_{n_1}}
\cong
\underbrace{\frac{m\mathbb Z}{u\mathbb Z}}_{\cong\,\mathbb Z_{u_1}}.
\]
% Indicii n_1 si u_1 de la baza paginii sunt usor estompati.
\[
u=m u_1,\qquad n=d n_1.
\]
\[
u_1=n_1\quad\Longrightarrow\quad
\frac{u}{m}=\frac{n}{d}\quad\Longrightarrow\quad du=mn.
\]
.
% Fara preambul si fara intrerupere fortata de pagina.
% Transcriere fidela; incertitudinile sunt mentionate in comentarii.

\subsection*{Curs 11 - Pagina 6}

% Continuarea demonstratiei teoremei de izomorfism, de la pagina 5.
\noindent $\underbrace{\bar a}_{aH}\in\ker f^*
\iff aH=aN\iff a\in H$.
% In primul rand al manuscrisului, aH=aN si a\in H sunt citiri
% probabile, dar compatibilitatea lor cu morfismul nu este clara.
\[
\ker f^*=\{x\mid (\ldots)\} = \frac{H}{N}.
\]
% Textul dintre bara verticala si acolada este partial barat/ilizibil.
% De asemenea, H/N este ceea ce pare scris; a nu se interpreta
% ca o corectie matematica a nucleului morfismului.

\bigskip
\noindent\underline{Aplica\c{t}ii:} \quad $(\mathbb Z,+,0)$.

\medskip
\noindent $H\leq\mathbb Z\ \Rightarrow\ (\exists!)\,u\in\mathbb N^*$
 astfel \^{i}nc\^{a}t $H=u\mathbb Z$.
% Manuscrisul pare sa foloseasca N^*; cazul H={0} cere u=0,
% deci enuntul, luat literal, nu acopera subgrupul trivial.

\bigskip
\noindent Fie $m,n\in\mathbb N^*$. Ar\u{a}t c\u{a}
\[
mn=[m,n]\cdot(m,n).
\]
\noindent % Notatii din manuscris: [m,n] = c.m.m.m.c.; (m,n) = c.m.m.d.c.
\[
m\mathbb Z+n\mathbb Z\leq\mathbb Z
\ \Longrightarrow\ (\exists!)\,d\in\mathbb N\ \text{a.\^{i}.}\
m\mathbb Z+n\mathbb Z=d\mathbb Z.
\]
\noindent Ar\u{a}t c\u{a} $d=(m,n)$.

\medskip
\noindent $(\exists)\,x,y\in\mathbb Z$ a.\^{i}.
\[
mx+ny=d
\quad\Longrightarrow\quad
\text{dac\u{a} } c\mid m\ \text{\c{s}i}\ c\mid n,
\ \text{atunci}\ c\mid d.
\]
\[
m\mathbb Z\subseteq d\mathbb Z,\qquad
m\in m\mathbb Z\ \Longrightarrow\ (\exists)\,q\in\mathbb Z
\ \text{a.\^{i}.}\ m=dq,\quad d\mid m.
\]
\noindent Analog $d\mid n$ $\Longrightarrow d=(m,n)$.

\bigskip
\noindent $m\mathbb Z\cap n\mathbb Z\leq\mathbb Z\ \Longrightarrow\
(\exists!)\,u\in\mathbb N$ a.\^{i}.
\[
m\mathbb Z\cap n\mathbb Z=u\mathbb Z.
\]
\noindent $u\in m\mathbb Z$ \c{s}i $u\in n\mathbb Z$
$\Longrightarrow m\mid u$ \c{s}i $n\mid u$.

\noindent Dac\u{a} $m\mid b$ \c{s}i $n\mid b$
$\stackrel{?}{\Longrightarrow} b\mid u$.
% In manuscris relatia finala pare b|u, dar randul urmator
% (b\in mZ si b\in nZ => b\in uZ) implica u|b. Verificare necesara.
\[
\underbrace{b\in m\mathbb Z}_{m\mid b},\quad
\underbrace{b\in n\mathbb Z}_{n\mid b}
\quad\Longrightarrow\quad b\in u\mathbb Z.
\]

\bigskip
\noindent T I de iz. $\Longrightarrow$
\[
\frac{m\mathbb Z+n\mathbb Z}{n\mathbb Z}
\cong \frac{m\mathbb Z}{m\mathbb Z\cap n\mathbb Z}.
\]
\[
\Longrightarrow\qquad
\underbrace{\frac{d\mathbb Z}{n\mathbb Z}}_{\cong\,\mathbb Z_{n_1}}
\cong
\underbrace{\frac{m\mathbb Z}{u\mathbb Z}}_{\cong\,\mathbb Z_{u_1}}.
\]
% Indicii n_1 si u_1 de la baza paginii sunt usor estompati.
\[
u=m u_1,\qquad n=d n_1.
\]
\[
u_1=n_1\quad\Longrightarrow\quad
\frac{u}{m}=\frac{n}{d}\quad\Longrightarrow\quad du=mn.
\]
.
% Fara preambul si fara intrerupere fortata de pagina.
% Transcriere fidela; incertitudinile sunt mentionate in comentarii.

\subsection*{Curs 11 - Pagina 6}

% Continuarea demonstratiei teoremei de izomorfism, de la pagina 5.
\noindent $\underbrace{\bar a}_{aH}\in\ker f^*
\iff aH=aN\iff a\in H$.
% In primul rand al manuscrisului, aH=aN si a\in H sunt citiri
% probabile, dar compatibilitatea lor cu morfismul nu este clara.
\[
\ker f^*=\{x\mid (\ldots)\} = \frac{H}{N}.
\]
% Textul dintre bara verticala si acolada este partial barat/ilizibil.
% De asemenea, H/N este ceea ce pare scris; a nu se interpreta
% ca o corectie matematica a nucleului morfismului.

\bigskip
\noindent\underline{Aplica\c{t}ii:} \quad $(\mathbb Z,+,0)$.

\medskip
\noindent $H\leq\mathbb Z\ \Rightarrow\ (\exists!)\,u\in\mathbb N^*$
 astfel \^{i}nc\^{a}t $H=u\mathbb Z$.
% Manuscrisul pare sa foloseasca N^*; cazul H={0} cere u=0,
% deci enuntul, luat literal, nu acopera subgrupul trivial.

\bigskip
\noindent Fie $m,n\in\mathbb N^*$. Ar\u{a}t c\u{a}
\[
mn=[m,n]\cdot(m,n).
\]
\noindent % Notatii din manuscris: [m,n] = c.m.m.m.c.; (m,n) = c.m.m.d.c.
\[
m\mathbb Z+n\mathbb Z\leq\mathbb Z
\ \Longrightarrow\ (\exists!)\,d\in\mathbb N\ \text{a.\^{i}.}\
m\mathbb Z+n\mathbb Z=d\mathbb Z.
\]
\noindent Ar\u{a}t c\u{a} $d=(m,n)$.

\medskip
\noindent $(\exists)\,x,y\in\mathbb Z$ a.\^{i}.
\[
mx+ny=d
\quad\Longrightarrow\quad
\text{dac\u{a} } c\mid m\ \text{\c{s}i}\ c\mid n,
\ \text{atunci}\ c\mid d.
\]
\[
m\mathbb Z\subseteq d\mathbb Z,\qquad
m\in m\mathbb Z\ \Longrightarrow\ (\exists)\,q\in\mathbb Z
\ \text{a.\^{i}.}\ m=dq,\quad d\mid m.
\]
\noindent Analog $d\mid n$ $\Longrightarrow d=(m,n)$.

\bigskip
\noindent $m\mathbb Z\cap n\mathbb Z\leq\mathbb Z\ \Longrightarrow\
(\exists!)\,u\in\mathbb N$ a.\^{i}.
\[
m\mathbb Z\cap n\mathbb Z=u\mathbb Z.
\]
\noindent $u\in m\mathbb Z$ \c{s}i $u\in n\mathbb Z$
$\Longrightarrow m\mid u$ \c{s}i $n\mid u$.

\noindent Dac\u{a} $m\mid b$ \c{s}i $n\mid b$
$\stackrel{?}{\Longrightarrow} b\mid u$.
% In manuscris relatia finala pare b|u, dar randul urmator
% (b\in mZ si b\in nZ => b\in uZ) implica u|b. Verificare necesara.
\[
\underbrace{b\in m\mathbb Z}_{m\mid b},\quad
\underbrace{b\in n\mathbb Z}_{n\mid b}
\quad\Longrightarrow\quad b\in u\mathbb Z.
\]

\bigskip
\noindent T I de iz. $\Longrightarrow$
\[
\frac{m\mathbb Z+n\mathbb Z}{n\mathbb Z}
\cong \frac{m\mathbb Z}{m\mathbb Z\cap n\mathbb Z}.
\]
\[
\Longrightarrow\qquad
\underbrace{\frac{d\mathbb Z}{n\mathbb Z}}_{\cong\,\mathbb Z_{n_1}}
\cong
\underbrace{\frac{m\mathbb Z}{u\mathbb Z}}_{\cong\,\mathbb Z_{u_1}}.
\]
% Indicii n_1 si u_1 de la baza paginii sunt usor estompati.
\[
u=m u_1,\qquad n=d n_1.
\]
\[
u_1=n_1\quad\Longrightarrow\quad
\frac{u}{m}=\frac{n}{d}\quad\Longrightarrow\quad du=mn.
\]
